% ===================================================================
%  TAULUKKO N:o 7 — Kylä talvella. — Maatalo keväällä.
%  Traduction 2026 d'apres texto/io, controlee sur texto/fr et
%  texto/sv. Consigne : voir texto/fi/10-taulukko-01.tex.
%
%  DEUX INVERSIONS EVITEES D'AVANCE, et toutes deux par la relative.
%  Au bloc c1-10-1, « la roto (17) glaciizesis dal glacio (25) di la
%  rivereto (23) » : le genitif finnois preposerait la riviere devant
%  la glace, et l'on ecrit donc « jaahan (25), joka on purossa (23) ».
%  Au bloc c2-08-1, « la kabano (38) di la hundo (37) » : meme piege,
%  meme remede — « kopin (38) vieresta, jossa koira (37) asuu ». Le
%  suedois avait bute au meme endroit et s'en etait tire de la meme
%  facon.
% ===================================================================

%%K t07-noto-1 noto
\VUnotes{143.2mm}{%
(*) Katso aterioiden yksityiskohdista taulukot 6 ja 13.%
}

%%K t07-apar-1 apar
\VUsaut{4.0mm}
\VUcentre{13.2pt}{90}{TOINEN SARJA}
\VUsaut{3.0mm}
\VUfilet{16mm}
\VUsaut{5.6mm}
\VUtitre{15.0pt}{60}{TAULUKKO N:o 7}
\VUsaut{3.0mm}

%%K t07-tit-1 sub
\VUcentre{12.0pt}{0}{\textit{Ensimmäinen kohtaus.}}
\VUsaut{3.0mm}

%%K t07-tit-2 sub
\VUpk{10.2pt}{40}{Kylä talvella.}
\VUsaut{4.0mm}

%%K t07-c1-01-1 p
1. --- Uudenvuoden loman aikana seurasin isääni Uuteenkaupunkiin, pieneen
kylään, jossa hänen oli hoidettava muutamia asioita.

%%K t07-c1-02-1 p
2. --- Otimme \VUgras{postivaunut} \textsuperscript{(11)}, jotka kulkevat
kaupungin ja tuon kauppalan väliä; mutta koska oli hyvin kylmä, tuo matka
vähän mukavassa ajopelissä ei ollut kovin miellyttävä.

%%K t07-c1-03-1 p
3. --- Siksi tunsimme todellista mielihyvää, kun näimme \VUgras{ajurin}
pysäyttävän vaununsa torille, \textit{Valkoisen Karhun} \VUgras{majatalon}
\textsuperscript{(2)} eteen. \VUgras{Majatalon isäntä} \textsuperscript{(4)}
odotti meitä talonsa kynnyksellä, samalla kun hän rauhallisesti poltti
\VUgras{piippuaan}.

%%K t07-c1-03-2 p
Hän pyysi meitä astumaan sisään, enkä minä antanut pyytää itseäni istumaan
risukasan tulen eteen, joka rätisi tulisijassa. Silloin pilkkasin purevaa
\VUgras{pohjoistuulta}, joka saa vanhan \VUgras{kyltin} \textsuperscript{(3)}
narisemaan ja joka mustina pyörteinä vei mukanaan \VUgras{savun}
\textsuperscript{(1)}, joka nousi savupiipusta.

%%K t07-c1-04-1 p
4. --- Oli aamiaisen aika. Odottamatta kauemmin söimme hyvällä ruokahalulla sen
yksinkertaisen mutta maukkaan aterian, joka meille tarjoiltiin (*).

%%K t07-tit-3 sub
\VUpk{10.2pt}{40}{Muutamia käyntejä.}
\VUsaut{3.0mm}

%%K t07-c1-05-1 p
5. --- Aamiaisen jälkeen seurasin isääni, joka on vakuutusasiamies, hänen
asiakkaidensa luo. Ensin kävimme \VUgras{sepän} \textsuperscript{(5)} luona,
joka asuu majatalon vieressä. Hän oli kengittämässä hevosta. Ihailin hänen
toverinsa voimaa, joka piti tukevasti hevosen \VUgras{kaviota}
nahkaesiliinallaan, samalla kun seppä kiinnitti \VUgras{kengän}
\textsuperscript{(6)} voimakkain vasarniskuin.

%%K t07-c1-06-1 p
6. --- Sitten menimme kylän \VUgras{suutarin} \textsuperscript{(7)} luo, vanhan
Jaakobin. Vaikka hänellä on komea \VUgras{saapas} kylttinään, hän tyytyi
puolipohjaamaan vanhan parin puhkikuluneita kenkiä.

%%K t07-c1-06-2 p
Vanhus sanoi meille nauraen: « Kaupungissa olin “suutari” eikä minulla ollut
asiakkaita. Täällä olen “kengänpaikkaaja” ja työtä on paljon. »

%%K t07-c1-07-1 p
7. --- Hän puhui meille naapuristaan \VUgras{parturista} \textsuperscript{(8)},
jonka asiakkaat eivät ole tyytyväisiä. Hänen \VUgras{partaveitsensä} ovat aina
lovilla ja hänen \VUgras{saksensa} eivät koskaan leikkaa hyvin. Mutta koska hän
on kylän ainoa parturi, on tietenkin pakko antaa hänen ajaa partansa ja
leikata hiuksensa. Muutoinkin hän on saita ja epäystävällinen mies.

%%K t07-c1-08-1 p
8. --- Onneksi muut \VUgras{naapurit} eivät ole hänen kaltaisiaan. Esimerkiksi
\VUgras{vaunuseppä} \textsuperscript{(9)} on hyväsydäminen mies, ahkera ja
avulias. On mielihyvä antaa hänen korjata vaunut. Hänen työnsä on aina
viimeistelty.

%%K t07-c1-09-1 p
9. --- Vanha Jaakob ei valittanut myöskään \VUgras{satulasepästä}
\textsuperscript{(10)}, jota hän joskus auttaa ompelemaan \VUgras{valjaita},
kun työllä on kiire.

%%K t07-c1-09-2 p
Isäni kysyi häneltä hänen perheestään: « Kaikki voivat hyvin. Lapsenlapseni käy
\VUgras{koulua} \textsuperscript{(12)}. Hänen \VUgras{opettajattarensa}
\textsuperscript{(13)} on hyvin tyytyväinen häneen, hän on hyvä \VUgras{oppilas}
\textsuperscript{(14)}. Hän ei ole vintiöpoikani kaltainen; tämä ajattelee vain
leikkimistä. \VUgras{Opettaja} \textsuperscript{(15)} ei onnistu opettamaan
hänelle yhtään mitään. »

%%K t07-tit-4 sub
\VUsaut{3.0mm}
\VUpk{10.2pt}{40}{Kyläjuoruja.}
\VUsaut{3.0mm}

%%K t07-c1-10-1 p
10. --- Vanhus kertoi meille sitten kylän uutiset. On määrä rakentaa uusi koulu,
sillä se, joka on sijoitettu \VUgras{kunnantaloon}, on käynyt liian pieneksi.
Muutoin se on helppoa, sillä riittää ostaa osa \VUgras{myllärin} puutarhasta. Se
on hyvin tuottoisaa poloiselle miehelle. Kylmyyden vuoksi \VUgras{myllyn}
\textsuperscript{(16)} \VUgras{myllynkivet} eivät enää voi jauhaa viljaa, sillä
\VUgras{pyörä} \textsuperscript{(17)} on jäätynyt kiinni \VUgras{jäähän}
\textsuperscript{(25)}, joka on \VUgras{purossa} \textsuperscript{(23)}.

%%K t07-c1-11-1 p
11. --- « Rakennuksista puheen ollen, oletteko nähneet uuden \VUgras{kirkkomme}
\textsuperscript{(18)}? Se on hyvin maalauksellinen lumivaippansa alla.
\VUgras{Varikset} \textsuperscript{(19)}, jotka eivät enää tunne vanhaa
kellotorniaan, istuvat surullisina suuressa puussa \VUgras{hautausmaalla}
\textsuperscript{(20)} ».

%%K t07-c1-12-1 p
12. --- Tuo ajatus hautausmaasta muistutti suutaria niistä henkilöistä, jotka
isäni oli tuntenut ja jotka kuolivat viime vuonna. « Kuinka monta \VUgras{hautaa}
\textsuperscript{(21)}, hyvä herra, onkaan lisätty toisten viereen
\VUgras{sypressin} \textsuperscript{(22)} alle! Kuolema on niittänyt vanhan
notaarimme. Kaikki kyläläiset olivat läsnä hänen \VUgras{hautajaisissaan}. »

%%K t07-c1-13-1 p
13. --- « Tässä on hänen seuraajansa, joka luistelee purolla. Hän tuli juuri
luokseni korjauttamaan \VUgras{luistimensa} \textsuperscript{(26)} hihnan, joka
oli katkennut. »

%%K t07-c1-14-1 p
14. --- « Tässä on myös puron rannalla uusi \VUgras{nimismies}
\textsuperscript{(27)}. Hän on entinen santarmi, joka pelottaa kylän vintiöitä.
He pakenevat heti kun saavat näkyviinsä hänen \VUgras{säärystimensä}
\textsuperscript{(29)} ja \VUgras{lippalakkinsa} \textsuperscript{(28)}. »

%%K t07-c1-15-1 p
15. --- « Ha! tässä on vanha Joosef, joka ajaa \VUgras{risukärryjä}
\textsuperscript{(30)} leipurille. --- Se on totta, sanoi isäni, tunnen hänen
\VUgras{muulivaljakkonsa} \textsuperscript{(31)}. Minun täytyy todellakin puhua
hänen kanssaan, käytän tilaisuuden hyväkseni. Hyvästi, isä Jaakob. »

%%K t07-c1-16-1 p
16. --- Juuri kun menimme ulos, kylän lapset lähtivät koulusta ja syöksyivät
juosten \VUgras{luistinradalle} \textsuperscript{(33)} vaarassa kaatua ja
katkaista raaja \VUgras{kaatuessaan} \textsuperscript{(34)}. Yksi heistä otti
\VUgras{kelkkansa} \textsuperscript{(32)} vaunusepältä, pani yhden tovereistaan
sille ja lähti juosten matkaan.

%%K t07-c1-17-1 p
17. --- Toiset syöksyivät \VUgras{lumikinokselle} \textsuperscript{(37)}
\VUgras{sillan} \textsuperscript{(40)} vieressä ja alkoivat pommittaa
\VUgras{lumiukkoa} \textsuperscript{(39)}, jonka he olivat pystyttäneet edellisenä
päivänä ja jolle he olivat panneet hatun, piipun ja koululaukun olalle.

%%K t07-c1-18-1 p
18. --- He välittävät vähän jäätävästä tuulesta ja kylmyydestä. Useimmilla ei ole
edes \VUgras{paksua kaulahuivia} \textsuperscript{(35)} ja, lämmittääkseen
itseään \VUgras{lumipallo}-taistelun \textsuperscript{(38)} jälkeen, heille
riittää kourallinen hyvin kuumia \VUgras{kastanjoita} \textsuperscript{(42)},
ostettuina vanhalta \VUgras{myyjättäreltä} Jaakobiinalta, joka värisee vilusta
liian ohuen \VUgras{repaleisen huivinsa} \textsuperscript{(43)} alla.

%%K t07-tit-5 sub
\VUsaut{3.0mm}
\VUcentre{12.0pt}{0}{\textit{Toinen kohtaus.}}
\VUsaut{3.0mm}

%%K t07-tit-6 sub
\VUpk{10.2pt}{40}{Maatalo keväällä.}
\VUsaut{4.0mm}

%%K t07-c2-01-1 p
1. --- Kaikista talven huveista huolimatta tervehdimme syvällä ilolla
\VUgras{kevään} paluuta.

%%K t07-c2-01-2 p
Minkä elävöittävän taulun tarjoaakaan maatila, illalla, toukokuussa!

%%K t07-c2-02-1 p
2. --- Taivaanrannalla \VUgras{aurinko} \textsuperscript{(1)} laskee hitaasti ja
tulvii \VUgras{maaseudun} \textsuperscript{(3)} purppuraisilla
\VUgras{säteillään} \textsuperscript{(2)}. Ahkera \VUgras{kyntäjä}
\textsuperscript{(6)} kiiruhtaa kyntämään \VUgras{peltonsa}
\textsuperscript{(4)}.

%%K t07-c2-02-2 p
\VUgras{Auransa} \textsuperscript{(5)} kanssa, jonka eteen on valjastettu
voimakas \VUgras{hevonen} \textsuperscript{(7)}, hän vetää \VUgras{vaon}
\textsuperscript{(8)}, johon \VUgras{kylväjä} \textsuperscript{(9)} täydellä
kädellä heittää \VUgras{siemenen}, joka antaa kultaiset tähkät.

%%K t07-c2-03-1 p
3. --- \VUgras{Niittomiehet} \textsuperscript{(11)} viikatteineen ovat niittäneet
niityn ruohon, heinänkääntäjät ovat kuivanneet \VUgras{heinän}
\textsuperscript{(10)} kääntämällä sitä \VUgras{heinähangoilla}
\textsuperscript{(12)} ja haravoilla, ja tässä he palaavat.

%%K t07-c2-03-2 p
Muutamat kulkevat raskaiden kärryjen edellä, joita härät vetävät; he kantavat
työkaluaan olallaan. Toiset, laiskemmat, ovat mukavasti asettuneet tuoksuvalle
heinälle.

%%K t07-c2-04-1 p
4. --- Ohi kulkiessaan he tervehtivät ystävällisesti \VUgras{puutarhuria}
\textsuperscript{(16)}, joka \VUgras{hedelmätarhassa} \textsuperscript{(13)} on
askaroimassa leikkaamassa \VUgras{hedelmäpuita} ja oksastamassa
\VUgras{päärynäpuita} \textsuperscript{(14)}. \VUgras{Luumupuut}
\textsuperscript{(15)} alkavat kukkia, ja hyvin lannoitetussa
\VUgras{kasvitarhassa} \textsuperscript{(17)} vihannekset, kaali ja salaatti
menestyvät jo hyvin.

%%K t07-c2-04-2 p
Pienellä muurilla kaunis \VUgras{riikinkukko} \textsuperscript{(18)} levittää
pyrstönsä antaakseen ihailla komeita höyheniään.

%%K t07-tit-7 sub
\VUpk{10.2pt}{40}{Maatila.}
\VUsaut{3.0mm}

%%K t07-c2-05-1 p
5. --- \VUgras{Tallin} \textsuperscript{(19)} ovet ovat auki ja näkyy seimi ja
\VUgras{olkivuode} \textsuperscript{(23)} \VUgras{tammalle}
\textsuperscript{(21)}, jonka \VUgras{tallirenki} \textsuperscript{(20)} on
juuri riisunut valjaista. \VUgras{Varsa} \textsuperscript{(22)}, niin
koomillinen pitkine jalkoineen, seuraa emoaan hyppelehtien.

%%K t07-c2-06-1 p
6. --- \VUgras{Kaivolle} \textsuperscript{(29)} menevät \VUgras{lehmät}
\textsuperscript{(25)} juomaan puisesta \VUgras{kaukalosta}
\textsuperscript{(26)}, joka on niiden juoma-altaana. \VUgras{Vasikka}
\textsuperscript{(24)} käyttää tilaisuutta imeäkseen, ja \VUgras{sonni}
\textsuperscript{(27)}, joka haistaa heinäladon hyvän tuoksun, ammuu iloisesti
ja kohottaa ylpeästi päätään voimakkaine \VUgras{sarvineen}
\textsuperscript{(28)}.

%%K t07-c2-07-1 p
7. --- Kaikki työskentelevät. \VUgras{Palvelustyttö} \textsuperscript{(32)}
pitää kädessään kaivon \VUgras{köyttä} \textsuperscript{(31)}. Hän ammentaa vettä
\VUgras{sangolla} \textsuperscript{(30)} juottaakseen karjan. \VUgras{Ladon}
\textsuperscript{(33)} vieressä mies haravoi olkea kokoon, ja
\VUgras{asuinrakennuksen} \textsuperscript{(34)} oven edessä vanha
\VUgras{kehrääjätär} \textsuperscript{(35)} kehrää rukkinsa ja \VUgras{värttinänsä}
kanssa \VUgras{pellavaa} tai \VUgras{hamppua} kuten muinaisina aikoina.

%%K t07-tit-8 sub
\VUsaut{3.0mm}
\VUpk{10.2pt}{40}{Maanviljelijä.}
\VUsaut{3.0mm}

%%K t07-c2-08-1 p
8. --- \VUgras{Maanviljelijä} \textsuperscript{(36)} johtaa kaikkia noita töitä
ja antaa ohjeita rengeilleen. Hän käskee viedä katon alle \VUgras{äkeen}
\textsuperscript{(40)} ja \VUgras{kuokan} \textsuperscript{(39)}, jotka on
unohdettu \VUgras{kopin} \textsuperscript{(38)} viereen, jossa \VUgras{koira}
\textsuperscript{(37)} asuu.

%%K t07-c2-09-1 p
9. --- Hänen huutonsa säikäyttää \VUgras{kyyhkysen} \textsuperscript{(41)}, joka
kaikin siivin lentää \VUgras{kyyhkyslakkaan} \textsuperscript{(42)}, ja
\VUgras{kanat} \textsuperscript{(43)}, jotka kiiruhtavat menemään takaisin
\VUgras{kanalaan} \textsuperscript{(44)}.

%%K t07-c2-10-1 p
10. --- \VUgras{Lammastarhan} \textsuperscript{(48)} edessä, tässä on vanha
\VUgras{paimen} \textsuperscript{(45)}, joka tuo takaisin \VUgras{(lammas)laumansa}
\textsuperscript{(49)}. \VUgras{Paimensauva} \textsuperscript{(46)} kädessään,
joka ei koskaan puutu häneltä, yhtä vähän kuin hänen haalistunut
\VUgras{pusakkansa} \textsuperscript{(47)}, hän ajaa tarhaan \VUgras{pässit}
\textsuperscript{(50)}, \VUgras{uuhet} \textsuperscript{(53)}, \VUgras{karitsat}
\textsuperscript{(54)}, \VUgras{vuohet} \textsuperscript{(51)} ja
\VUgras{kilit} \textsuperscript{(52)}. Hänen uskollinen \VUgras{koiransa}
\textsuperscript{(55)} Argus tulee viimeisenä ja yllyttää haukunnallaan
jälkeenjääneiden joutuisuutta.

%%K t07-c2-11-1 p
11. --- Kaikki tuo melu ja liike ei häiritse hyvän \VUgras{lypsylehmän}
\textsuperscript{(56)} rauhaa, joka antaa \VUgras{kellonsa}
\textsuperscript{(57)} kilistä, samalla kun sitä lypsetään \VUgras{navetan}
\textsuperscript{(58)} oven edessä.

%%K t07-c2-12-1 p
12. Se näyttää ymmärtävän, että sen täytyy antaa maitonsa, jotta
\VUgras{emäntä} \textsuperscript{(60)} voisi tehdä \VUgras{voita}
\VUgras{kirnussa} \textsuperscript{(61)} tai erinomaisia \VUgras{juustoja}
\textsuperscript{(64)}, jotka tippuvat hyllyltä, joka on \VUgras{saavin}
\textsuperscript{(63)} päällä.

%%K t07-c2-13-1 p
13. --- Koko \VUgras{maitohuoneen} \textsuperscript{(59)} pitää hyvin puhtaana
\VUgras{pihamies} \textsuperscript{(65)}, joka on juuri pessyt
\VUgras{maitoruukut} \textsuperscript{(62)} siinä suuressa sangossa, jota hän
kantaa kädessään.

%%K t07-tit-9 sub
\VUsaut{3.0mm}
\VUpk{10.2pt}{40}{Siipikarja.}
\VUsaut{3.0mm}

%%K t07-c2-14-1 p
14. --- Paksuine puukenkineen hän kävelee hieman raskaasti, ja niin hän
karkottaa \VUgras{kalkkunan} \textsuperscript{(68)}, \VUgras{ankat}
\textsuperscript{(66)}, \VUgras{urosankat} \textsuperscript{(67)} ja
\VUgras{hanhet} \textsuperscript{(69)}, jotka ammottavat \VUgras{nokkaansa}
\textsuperscript{(70)} ja kirkuvat levottomasti.

%%K t07-c2-14-2 p
\VUgras{Kukko} \textsuperscript{(71)}, taistelunhaluisempi, kohottaa punaisen
\VUgras{heltansa} \textsuperscript{(72)}, ravistaa \VUgras{pyrstönsä}
\textsuperscript{(73)} höyheniä ja antaa kuulua äänekkään « kukkokiekuu ».

%%K t07-c2-15-1 p
15. --- \VUgras{Emokana} \textsuperscript{(74)} etsii ruokaa
\VUgras{lantakasalla} \textsuperscript{(81)} \VUgras{pienten tipujensa}
\textsuperscript{(75)} kanssa. \VUgras{Porsas} \textsuperscript{(78)} pakenee
emonsa luo, sen valtavan \VUgras{emakon} \textsuperscript{(77)}, jonka näemme
sikalan vieressä.

%%K t07-c2-16-1 p
16. --- Häkeissään \VUgras{kanit} \textsuperscript{(79)} syövät ahnaasti hentoa
\VUgras{ruohoa} \textsuperscript{(80)}, jota maanviljelijän tytär kantaa niille.
Pikku Johanna pitää kovasti kaneistaan ja, joka päivä, hän tulee tapaamaan
pieniä ystäviään sen jälkeen kun hän on hakenut tuoreet \VUgras{munat}
\textsuperscript{(76)} kanalasta.
\VUsaut{6.0mm}
\VUfilet{22mm}
